
% ----------------------------------------------------------------------
%                   LATEX TEMPLATE FOR PhD THESIS
% ----------------------------------------------------------------------

% based on Harish Bhanderi's PhD/MPhil template, then Uni Cambridge
% http://www-h.eng.cam.ac.uk/help/tpl/textprocessing/ThesisStyle/
% corrected and extended in 2007 by Jakob Suckale, then MPI-CBG PhD programme
% and made available through OpenWetWare.org - the free biology wiki


%: Style file for Latex
% Most style definitions are in the external file PhDthesisPSnPDF.
% In this template package, it can be found in ./Latex/Classes/
\documentclass[twoside,11pt]{Latex/Classes/PhDthesisPSnPDF}


%: Macro file for Latex
% Macros help you summarise frequently repeated Latex commands.
% Here, they are placed in an external file /Latex/Macros/MacroFile1.tex
% An macro that you may use frequently is the figuremacro (see introduction.tex)
% % This file contains macros that can be called up from connected TeX files
% It helps to summarise repeated code, e.g. figure insertion (see below).

% insert a centered figure with caption and description
% parameters 1:filename, 2:title, 3:description and label
\newcommand{\figuremacro}[3]{
	\begin{figure}[htbp]
		\centering
		\includegraphics[width=1\textwidth]{#1}
		\caption[#2]{\textbf{#2} - #3}
		\label{#1}
	\end{figure}
}

% insert a centered figure with caption and description AND WIDTH
% parameters 1:filename, 2:title, 3:description and label, 4: textwidth
% textwidth 1 means as text, 0.5 means half the width of the text
\newcommand{\figuremacroW}[4]{
	\begin{figure}[htbp]
		\centering
		\includegraphics[width=#4\textwidth]{#1}
		\caption[#2]{\textbf{#2} - #3}
		\label{#1}
	\end{figure}
}

% inserts a figure with wrapped around text; only suitable for NARROW figs
% o is for outside on a double paged document; others: l, r, i(inside)
% text and figure will each be half of the document width
% note: long captions often crash with adjacent content; take care
% in general: above 2 macro produce more reliable layout
\newcommand{\figuremacroN}[3]{
	\begin{wrapfigure}{o}{0.5\textwidth}
		\centering
		\includegraphics[width=0.48\textwidth]{#1}
		\caption[#2]{{\small\textbf{#2} - #3}}
		\label{#1}
	\end{wrapfigure}
}

% predefined commands by Harish
\newcommand{\PdfPsText}[2]{
  \ifpdf
     #1
  \else
     #2
  \fi
}

\newcommand{\IncludeGraphicsH}[3]{
  \PdfPsText{\includegraphics[height=#2]{#1}}{\includegraphics[bb = #3, height=#2]{#1}}
}

\newcommand{\IncludeGraphicsW}[3]{
  \PdfPsText{\includegraphics[width=#2]{#1}}{\includegraphics[bb = #3, width=#2]{#1}}
}

\newcommand{\InsertFig}[3]{
  \begin{figure}[!htbp]
    \begin{center}
      \leavevmode
      #1
      \caption{#2}
      \label{#3}
    \end{center}
  \end{figure}
}


%%% Local Variables: 
%%% mode: latex
%%% TeX-master: "~/Documents/LaTeX/CUEDThesisPSnPDF/thesis"
%%% End: 

\usepackage[T1]{fontenc}
\usepackage{array}
\usepackage{pdfpages}

%\usepackage{graphics}
% or use the graphicx package for more complicated commands
%\usepackage{graphicx}

%: ----------------------------------------------------------------------
%:                  TITLE PAGE: name, degree,..
% ----------------------------------------------------------------------
\usepackage{graphicx}

      \textwidth 15cm
      \textheight 22cm
      \parindent 10pt
      \oddsidemargin 0.85cm
      \evensidemargin 0.37cm


\begin{document}

\thispagestyle{empty}

\begin{center}

Vrije Universiteit Amsterdam

\vspace{1mm}

\includegraphics[height=25mm]{0_frontmatter/figures/vu-griffioen.pdf}

\vspace{2cm}

{\Large Bachelor Thesis}

\vspace*{1.5cm}

\rule{.9\linewidth}{.6pt}\\[0.4cm]
{\huge \bfseries RDMA Key Value store:\\Scalability of RDMA transportation types \par}\vspace{0.4cm}
\rule{.9\linewidth}{.6pt}\\[1.5cm]

\vspace*{2mm}

{\Large
\begin{tabular}{l}
{\bf Author:} ~~Jan Erik (Jens) Troost ~~~~ (2656054)
\end{tabular}
}

\vspace*{2cm}

\begin{tabular}{ll}
{\it 1st supervisor:}   & ~~Animesh Trivedi \\
{\it daily supervisor:} & ~~supervisor name \\
{\it 2nd reader:}       & ~~supervisor name
\end{tabular}

\vspace*{2.5cm}

\textit{A thesis submitted in fulfillment of the requirements for\\the VU Bachelor of Science degree in Computer Science}

\vspace*{1.8cm}

\today\\[4cm] % Date

\end{center}

\newpage


% ----------------------------------------------------------------------
       
% turn of those nasty overfull and underfull hboxes
\hbadness=10000
\hfuzz=50pt


%: --------------------------------------------------------------
%:                  FRONT MATTER: dedications, abstract,..
% --------------------------------------------------------------


%\language{english}


% sets line spacing
\renewcommand\baselinestretch{1.2}
\baselineskip=18pt plus1pt


%: ----------------------- generate cover page ------------------------

%\begin{center}
%\textit{``I am the master of my fate, I am the captain of my soul'' \\ from {\em Invictus}, by William Ernest Henley}
%\end{center}

%: ----------------------- cover page back side ------------------------
% Your research institution may require reviewer names, etc.
% This cover back side is required by Dresden Med Fac; uncomment if needed.

\newpage
%\vspace{10mm}
%1. First Reader: Name Surname
%
%\vspace{10mm}
%2. Daily Supervisor: Name Surname
%
%\vspace{10mm}
%3. Second Reader: Name Surname
%
%\vspace{10mm}
%4. Industrial Supervisor: Name Surname
%
%\vspace{20mm}
%Day of the defense:

%\vspace{20mm}
%\hspace{70mm}Signature from head of PhD committee:



%: ----------------------- abstract ------------------------

% Your institution may have specific regulations if you need an abstract and where it is to be placed in the document. The default here is just after title.


% Thesis Abstract -----------------------------------------------------


%\begin{abstractslong}    %uncommenting this line, gives a different abstract heading
\begin{abstracts}        %this creates the heading for the abstract page

    Seemingly ever growing tech giants, such as Facebook, Amazon and Google, require fast, reliable and scalable key-value storage (KV-store) to serve product recommendations, user preferences, and advertisements.
    Past advancements of KV-stores have focussed on improving these underlying data structures.
    Remote direct memory access (RDMA) networks have increasingly become more popular in commercial and academic data centers.
    This relatively new technology offer lower latency and higher throughput performance compared to tradition sockets and network interface cards (NIC).
    Scalability performance of RDMA transportation protocols is less known, previous work focused on RDMA verb choice.

    This paper focuses on this gap, evaluating scalability performance of RDMA transportation types.
    These findings will be used to give recommendation for RDMA KV-store implementations.
    Macro-level benchmarks have been conducted, with realistic KV-store workloads.
    It has been found that RDMA offers a significant improvements compared to TCP.
    UD has shown to perform best, with 94.8\% throughput improvement over TCP, and 58.7\% improvement over RC.
    All the while offering consistent and gradually increasing latency up to 30 clients, reaching a maximum average latency of 64 $\mu$sec.
    Additionally, scalability issues of RC have been shown, and is not recommended in a KV-store application.

\end{abstracts}
%\end{abstractlongs}


% ---------------------------------------------------------------------- 


% The original template provides and abstractseparate environment, if your institution requires them to be separate. I think it's easier to print the abstract from the complete thesis by restricting printing to the relevant page.
% \begin{abstractseparate}
%   
% Thesis Abstract -----------------------------------------------------


%\begin{abstractslong}    %uncommenting this line, gives a different abstract heading
\begin{abstracts}        %this creates the heading for the abstract page

    Seemingly ever growing tech giants, such as Facebook, Amazon and Google, require fast, reliable and scalable key-value storage (KV-store) to serve product recommendations, user preferences, and advertisements.
    Past advancements of KV-stores have focussed on improving these underlying data structures.
    Remote direct memory access (RDMA) networks have increasingly become more popular in commercial and academic data centers.
    This relatively new technology offer lower latency and higher throughput performance compared to tradition sockets and network interface cards (NIC).
    Scalability performance of RDMA transportation protocols is less known, previous work focused on RDMA verb choice.

    This paper focuses on this gap, evaluating scalability performance of RDMA transportation types.
    These findings will be used to give recommendation for RDMA KV-store implementations.
    Macro-level benchmarks have been conducted, with realistic KV-store workloads.
    It has been found that RDMA offers a significant improvements compared to TCP.
    UD has shown to perform best, with 94.8\% throughput improvement over TCP, and 58.7\% improvement over RC.
    All the while offering consistent and gradually increasing latency up to 30 clients, reaching a maximum average latency of 64 $\mu$sec.
    Additionally, scalability issues of RC have been shown, and is not recommended in a KV-store application.

\end{abstracts}
%\end{abstractlongs}


% ---------------------------------------------------------------------- 

% \end{abstractseparate}


%: ----------------------- tie in front matter ------------------------

\frontmatter
% Thesis Dedication ---------------------------------------------------

%\begin{dedication} %this creates the heading for the dedication page

%To ...

%\end{dedication}

% ----------------------------------------------------------------------
% Thesis Acknowledgements ------------------------------------------------


%\begin{acknowledgementslong} %uncommenting this line, gives a different acknowledgements heading
%\begin{acknowledgements}      %this creates the heading for the acknowlegments


%\end{acknowledgements}
%\end{acknowledgmentslong}

% ------------------------------------------------------------------------





%: ----------------------- contents ------------------------

\setcounter{secnumdepth}{3} % organisational level that receives a numbers
\setcounter{tocdepth}{3}    % print table of contents for level 3
\tableofcontents            % print the table of contents
% levels are: 0 - chapter, 1 - section, 2 - subsection, 3 - subsection


%: ----------------------- list of figures/tables ------------------------

\listoffigures	% print list of figures

\listoftables  % print list of tables




%: ----------------------- glossary ------------------------

% Tie in external source file for definitions: /0_frontmatter/glossary.tex
% Glossary entries can also be defined in the main text. See glossary.tex
% 
%% this file is called up by thesis.tex
% content in this file will be fed into the main document

% Glossary entries are defined with the command \nomenclature{1}{2}
% 1 = Entry name, e.g. abbreviation; 2 = Explanation
% You can place all explanations in this separate file or declare them in the middle of the text. Either way they will be collected in the glossary.

% required to print nomenclature name to page header
%\markboth{\MakeUppercase{\nomname}}{\MakeUppercase{\nomname}}


% ----------------------- contents from here ------------------------


%\nomenclature{LSY}{ehbfuefebbfbjkjkebfjbfbfw} 
%\nomenclature{DEPC}{diethyl-pyro-carbonate; used to remove RNA-degrading enzymes (RNAases) from water and laboratory utensils}
%\nomenclature{DMSO}{dimethyl sulfoxide; organic solvent, readily passes through skin, cryoprotectant in cell culture}
%\nomenclature{EDTA}{Ethylene-diamine-tetraacetic acid; a chelating (two-pronged) molecule used to sequester most divalent (or trivalent) metal ions, such as calcium (Ca$^{2+}$) and magnesium (Mg$^{2+}$), copper (Cu$^{2+}$), or iron (Fe$^{2+}$ / Fe$^{3+}$)}



 

%\begin{multicols}{2} % \begin{multicols}{#columns}[header text][space]
%\begin{footnotesize} % scriptsize(7) < footnotesize(8) < small (9) < normal (10)

%\printnomenclature[1.5cm] % [] = distance between entry and description
%\label{nom} % target name for links to glossary

%\end{footnotesize}
%\end{multicols}



%: --------------------------------------------------------------
%:                  MAIN DOCUMENT SECTION
% --------------------------------------------------------------

% the main text starts here with the introduction, 1st chapter,...
\mainmatter

\renewcommand{\chaptername}{} % uncomment to print only "1" not "Chapter 1"


%: ----------------------- subdocuments ------------------------

% Parts of the thesis are included below. Rename the files as required.
% But take care that the paths match. You can also change the order of appearance by moving the include commands.


% this file is called up by thesis.tex
% content in this file will be fed into the main document

%: ----------------------- introduction file header -----------------------
\chapter{Introduction}

% the code below specifies where the figures are stored
%\ifpdf
%    \graphicspath{{1_introduction/figures/PNG/}{1_introduction/figures/PDF/}{1_introduction/figures/}}
%\else
%    \graphicspath{{1_introduction/figures/EPS/}{1_introduction/figures/}}
%\fi

% ----------------------------------------------------------------------
%: ----------------------- introduction content ----------------------- 
% ----------------------------------------------------------------------



%: ----------------------- HELP: latex document organisation
% the commands below help you to subdivide and organise your thesis
%    \chapter{}       = level 1, top level
%    \section{}       = level 2
%    \subsection{}    = level 3
%    \subsubsection{} = level 4
% note that everything after the percentage sign is hidden from output



\section{Context}
Seemingly ever growing tech giants, such as Facebook, Amazon and Google, require fast, reliable and scalable key-value storage (KV-store) to serve a variety of services: product recommendations, user preferences, and advertisements\cite{decandia2007dynamo,geambasu2010comet}.
Amazon's Dynamo is an example of such a highly used KV-store, which is used by Amazon to provide its core services, such as session storage and shopping cart\cite{decandia2007dynamo}.
At peak, these services handle tens of millions customers a second.
This requires Dynamo to process requests in the span of a few milliseconds, to provide a consistent and fast user experience.

Remote direct memory access (RDMA) networks have increasingly become more popular in data centers, due to their increase in throughput, decrease in latency, and lower cost of RDMA capapble network interface cards (RNIC)\cite{kalia2016design, chen2019scalable}.
RDMA and RNIC provides a lower latency networking interface compared to classical network cards, which is advantage for latency sensitive applications, such as a KV-store, deployed in data centers.


In a more general sense, a KV-store aims to store many small values, identified with a key.
KV-stores usually offer a simple interface using 3 commands: \textit{GET}, \textit{SET}, and \textit{DEL}.
With these commands, data can be retrieved, updated/added, or deleted.
Typically, the basis of a KV-store is formed by a hash-table, which pairs a unique key with a value.
Past advancements of KV-stores have focussed on improving these underlying data structures\cite{escriva2012hyperdex, lim2014mica}.
Cuckoo and Hopscotch hashing\cite{kalia2014using, FaRM} have presented an alternative of an improved hash table algorithm.
Nonetheless, with recent developments with RDMA networks, this has become a more attractive direction to improve KV-stores\cite{chen2019scalable}.

Remote direct memory access (RDMA) networks have increasingly become more popular in commercial and academic data centers, due to the increase in availability and decrease in cost for RDMA capable network interface cards (RNICs)\cite{kalia2016design, chen2019scalable}.
RDMA offers a more direct connection between two machines, by using direct memory acces (DMA), without the involvement of host CPU or OS.

\section{Problem statement}
Making efficient use of RDMA networks is a difficult task, and requires in-depth knowledge on the hardware constraints present in the RNICs\cite{kalia2016design, chen2019scalable}.
Additionally, in a higher level sense, there are design choices to be made, in regard to transportation types and so-called verbs.
RDMA networks can handle of various transportation types: reliable connection (RC), unreliable connection (UC) and unreliable datagrams (UD).
Additionally, verbs can be seen as network operations.
There are four main verbs: \textit{SEND}, \textit{RECV}, \textit{READ}, and \textit{WRITE}.
Each transportation type and verb come with advantages and disadvantages, which will be discussed further in this thesis.
However, impact of transportation types on scalability of an RDMA bases KV-store has not been explicitly examined.
Previous work has focused on researching verb choices, less on transportation type\cite{kalia2014using, kalia2016fasst, mitchell2013using}.

Scalability of their KV-store service is an important factor for tech giants such as Amazon.
With their growth in number of customers and concurrent clients, their KV-store servers are hit with large number of clients.



Scalability is an important factor for growing tech giants, as their increasing customer .
These require a low latency and scale to reach large number of clients, all the while being highly available\cite{decandia2007dynamo}, which RDMA could provide.

\section{Research Question}
This paper will explore to what extent RDMA transportation types affect the performance and multicore scalability of KV-store.
For this, this research will answer the following questions:

\begin{itemize}
    \item[\textbf{RQ1}] What is the multicore scalability of the RDMA transportation types: reliable connection (RC), unreliable connection (UC), and unreliable datagram (UD)?
    With the growing number of clients that
    \item[\textbf{RQ2}] What are the advantages and disadvantages of RDMA transportation types on an KV-store?
    This question aims to aid with design process of future RDMA KV stores.
    By applying known ramifications of RDMA to that in a KV store use case.
    With supporting results, recommendations can be given, and possible design issues can be foreseen.
\end{itemize}

\section{Research Methods}
In this thesis the network performance of KV-store, with varying network protocols, will be studied.
First an understanding of KV-stores is established.
Along with discussing known issues with traditional networking, RDMA will be introduced.
Experiments will be performed to measure overall throughput and latency, with increasing number of clients.
With this the possible scalability of each transport type will be shown.

A prototype\cite{iosup2019atlarge,hamming1998art,peffers2007design} and experimental\cite{jain1990art,heiser2010,ousterhout2018always} approach is taken for this thesis:
\begin{itemize}
    \item[\textbf{M1}] A KV-store prototype will be implemented, with a flexible network interfacing, to include both traditional socket and RDMA interfacing.
    This, and all other relevant project files, are open source and can be found on Github\cite{github}.
    \item[\textbf{M2}] To investigate the performance and scalability between the various networking types, a workload realistic\cite{atikoglu2012workload} macro-level benchmark will be designed.
\end{itemize}

With the focus of this thesis being on the network implementation, a trivial KV-store will be used.
This implementation does not offer strong performance and scalability, however these issues are minimized and kept consistent throughout experiments.

All measurements hereafter are recorded on the DAS-5 computing cluster.
This allows for a consistent working environment.
Further details on this is shown in table \ref{tab:das5} and discussed in section \ref{subsec:experimental-setup}

\section{Thesis Contributions}
This thesis presents scalability performance of RC and UD are shown and compared to the well established TCP protocol.
Additionally, recommendation as to which transportation protocol is best suited for an RDMA KV-store application.
These findings can be used to further examine possible optimizations or RDMA verb choices, or aid in future RDMA KV-store implementations.
The prototype and all project files can be found on Github\cite{github}.

\section{Plagiarism Declaration}
I herby declare that this thesis work, all data collected and findings are my own.

For this thesis, the basic structure of KV store from course Operating Systems, at the VU, has been used.
Along with basic RDMA functions from Dr. Trivedi's RDMA example\cite{rdma_example}.

Note: permission incoming for use of KV store structure from OS course.

\section{Thesis Structure}
Section \ref{ch:background} will provide the necessarily background knowledge of KV-store, linux sockets, and RDMA.
Next, section \ref{ch:design} will go over the design of the KV-store, networking interfacing, and benchmark.
In section \ref{ch:design} this design will be taken and described the implementation in a more technical prospective.
The benchmarking result will be presented and analyzed in section \ref{ch:evaluation}.
Section \ref{ch:related-work} compares findings with that of previously done work.
Future prospects will be given in section \ref{ch:future-improvements}.
Closing off the thesis, section \ref{ch:conclusion} will go over the conclusion and provide recommendations.

% this file is called up by thesis.tex
% content in this file will be fed into the main document

\chapter{This is another chapter} % top level followed by section, subsection


% ----------------------- paths to graphics ------------------------

% change according to folder and file names
\ifpdf
    \graphicspath{{7/figures/PNG/}{7/figures/PDF/}{7/figures/}}
\else
    \graphicspath{{7/figures/EPS/}{7/figures/}}
\fi


% ----------------------- contents from here ------------------------
% 

\dots



% ---------------------------------------------------------------------------
% ----------------------- end of thesis sub-document ------------------------
% ---------------------------------------------------------------------------
	
%\include{3_materials&methods/materials_methods}						
%\include{4_analysis&results/analysis&results}	

%\include{5_discussions/discussions}

%% this file is called up by thesis.tex
% content in this file will be fed into the main document

\chapter{Conclusion}\label{ch:conclusion} % top level followed by section, subsection


% ----------------------- paths to graphics ------------------------

% change according to folder and file names
\ifpdf
    \graphicspath{{7/figures/PNG/}{7/figures/PDF/}{7/figures/}}
\else
    \graphicspath{{7/figures/EPS/}{7/figures/}}
\fi


% ----------------------- contents from here ------------------------
%
RDMA has been shown to be a promising advancement for key values stores.
The low latency that RDMA brings with it, opens op new possibilities and improved performance.
RDMA also requires more design decision to be made, which transportation protocol to use, which optimizations to use, and which verbs to use.
In this thesis, the transportation protocols RC and UD have been compared, on performance, to TCP.

\textbf{RQ1:} In industry, KV stores require to be scalable and low latency.
Connection based transportation (RC and UC) put stress on RNIC when using large number of clients, hindering performance with respect to throughput and scalability.
RC has similar scalability to TCP, with a near 60 kilo-tasks/sec increase in throughput across all number of clients.
Best performing transportation type is UD, with close to 400 kilo-tasks/sec throughput and 64 $\mu$sec at 30 clients.
It has been shown that UD offers a 58.7\% improvement in throughput against RC, and 45.5\% improvement in latency.
For scalability UD is recommended to use, as this offers the best overall performance and sustains this up to 30 clients, and has the potential for optimizations to further improve performance.

\textbf{RQ2:} The advantages and disadvantages present with the different transportation types impact the KV store design choices.
In context of KV stores, UD is a compelling transportation type, however does not offer RDMA verbs such as \textit{READ} and \textit{WRITE}, which could offer for improved performance.
UD also allow for more optimizations, due to their one-to-many QP.
UD performance in this thesis have found to be limited.
Multiple UD QP's along with client grouping could continue scalability further.
Connection based transportation types are limited by one-to-one QP, although also have some room for improvements.


% ---------------------------------------------------------------------------
% ----------------------- end of thesis sub-document ------------------------
% ---------------------------------------------------------------------------
          
%% this file is called up by thesis.tex
% content in this file will be fed into the main document

%: ----------------------- name of chapter  -------------------------
\chapter*{Appendix} % top level followed by section, subsection


%: ----------------------- paths to graphics ------------------------

% change according to folder and file names
\ifpdf
    \graphicspath{{X/figures/PNG/}{X/figures/PDF/}{X/figures/}}
\else
    \graphicspath{{X/figures/EPS/}{X/figures/}}
\fi

%: ----------------------- contents from here ------------------------
\includepdf[pages=1-5]{CodesTable} 
\includepdf[pages=1-4]{InterviewQuestions}





% ---------------------------------------------------------------------------
%: ----------------------- end of thesis sub-document ------------------------
% ---------------------------------------------------------------------------


      
            
% --------------------------------------------------------------
%:                  BACK MATTER: appendices, refs,..
% --------------------------------------------------------------

% the back matter: appendix and references close the thesis


%: ----------------------- bibliography ------------------------

% The section below defines how references are listed and formatted
% The default below is 2 columns, small font, complete author names.
% Entries are also linked back to the page number in the text and to external URL if provided in the BibTex file.

% PhDbiblio-url2 = names small caps, title bold & hyperlinked, link to page 
%\begin{multicols}{2} % \begin{multicols}{ # columns}[ header text][ space]
%\begin{tiny} % tiny(5) < scriptsize(7) < footnotesize(8) < small (9)

\bibliographystyle{Latex/Classes/PhDbiblio-url2} % Title is link if provided
\renewcommand{\bibname}{References} % changes the header; default: Bibliography

\bibliography{7_references/references} % adjust this to fit your BibTex file

%\end{tiny}
%\end{multicols}



% --------------------------------------------------------------
% Various bibliography styles exit. Replace above style as desired.

% in-text refs: (1) (1; 2)
% ref list: alphabetical; author(s) in small caps; initials last name; page(s)
%\bibliographystyle{Latex/Classes/PhDbiblio-case} % title forced lower case
%\bibliographystyle{Latex/Classes/PhDbiblio-bold} % title as in bibtex but bold
%\bibliographystyle{Latex/Classes/PhDbiblio-url} % bold + www link if provided

%\bibliographystyle{Latex/Classes/jmb} % calls style file jmb.bst
% in-text refs: author (year) without brackets
% ref list: alphabetical; author(s) in normal font; last name, initials; page(s)

%\bibliographystyle{plainnat} % calls style file plainnat.bst
% in-text refs: author (year) without brackets
% (this works with package natbib)


% --------------------------------------------------------------

% according to Dresden med fac summary has to be at the end
%
% Thesis Abstract -----------------------------------------------------


%\begin{abstractslong}    %uncommenting this line, gives a different abstract heading
\begin{abstracts}        %this creates the heading for the abstract page

    Seemingly ever growing tech giants, such as Facebook, Amazon and Google, require fast, reliable and scalable key-value storage (KV-store) to serve product recommendations, user preferences, and advertisements.
    Past advancements of KV-stores have focussed on improving these underlying data structures.
    Remote direct memory access (RDMA) networks have increasingly become more popular in commercial and academic data centers.
    This relatively new technology offer lower latency and higher throughput performance compared to tradition sockets and network interface cards (NIC).
    Scalability performance of RDMA transportation protocols is less known, previous work focused on RDMA verb choice.

    This paper focuses on this gap, evaluating scalability performance of RDMA transportation types.
    These findings will be used to give recommendation for RDMA KV-store implementations.
    Macro-level benchmarks have been conducted, with realistic KV-store workloads.
    It has been found that RDMA offers a significant improvements compared to TCP.
    UD has shown to perform best, with 94.8\% throughput improvement over TCP, and 58.7\% improvement over RC.
    All the while offering consistent and gradually increasing latency up to 30 clients, reaching a maximum average latency of 64 $\mu$sec.
    Additionally, scalability issues of RC have been shown, and is not recommended in a KV-store application.

\end{abstracts}
%\end{abstractlongs}


% ---------------------------------------------------------------------- 


%: Declaration of originality
%\include{8_backmatter/declaration}



\end{document}
